\documentclass[10pt,a4paper]{article}
\usepackage[margin=0.85in]{geometry}
\usepackage[T1]{fontenc}
\usepackage[utf8]{inputenc}
\usepackage{lmodern,amsmath,amssymb,graphicx,enumitem,microtype,fancyhdr,xcolor,needspace}
\usepackage[hidelinks]{hyperref}
\definecolor{lightgrey}{gray}{0.94}
\definecolor{midgrey}{gray}{0.35}
\setlength{\parindent}{0pt}
\setlength{\parskip}{5pt}
\setlist[enumerate]{itemsep=5pt,topsep=4pt}
\pagestyle{fancy}\fancyhf{}\renewcommand{\headrulewidth}{0pt}
\fancyfoot[L]{\footnotesize\color{midgrey}Arij Asad}
\fancyfoot[C]{\footnotesize\color{midgrey}Edexcel companion practice}
\fancyfoot[R]{\footnotesize\color{midgrey}\thepage}
\newcounter{question}
\newenvironment{question}[1]{\refstepcounter{question}\Needspace{5\baselineskip}\par\medskip\textbf{\thequestion.\ #1}\par}{\par\vspace{12mm}}
\begin{document}
\hypersetup{pdftitle={Pure Mathematics Year 2 - Original practice and solutions},pdfauthor={Arij Asad}}
\begin{minipage}[c]{0.78\textwidth}{\Large\bfseries Pure Mathematics Year 2}\par Original practice check\end{minipage}\hfill\begin{minipage}[c]{0.17\textwidth}\IfFileExists{PS_logo.png}{\includegraphics[width=\linewidth]{PS_logo.png}}{\rule{0pt}{16mm}}\end{minipage}

\medskip\fcolorbox{black!20}{lightgrey}{\parbox{\dimexpr\linewidth-2\fboxsep-2\fboxrule\relax}{Three short checks on trigonometric equations, parametric tangents and integration by parts. Angles are in radians. Attempt the questions before reading the solutions. These original questions sample selected skills; they are not an official Edexcel paper.}}

\begin{question}{Trigonometric equations}
Solve \(2\sin^2x+\sin x-1=0\) for \(0\leq x<2\pi\).
\end{question}
\begin{question}{Parametric differentiation}
A curve is given by \(x=t^2+1\), \(y=t^3-3t\). Find the equation of its tangent at \(t=2\).
\end{question}
\begin{question}{Integration by parts}
Evaluate \(\displaystyle\int_0^1 xe^x\,\mathrm{d}x\).
\end{question}
\Needspace{7\baselineskip}\textbf{Find the mistake}\par
A student calculates \(\int_{-1}^1x\,\mathrm{d}x=0\) and concludes that the total area between \(y=x\) and the horizontal axis on this interval is zero. Correct the conclusion.

\vfill\small Further practice: \href{https://maths.arijasad.com/tmua-paper-archive.html}{TMUA papers and solutions} and \href{https://maths.arijasad.com/maths-topic-practice.html}{maths topic guides}.
\newpage
{\Large\bfseries Hints and worked solutions}\par\medskip
\Needspace{8\baselineskip}\subsection*{1. Trigonometric equations}
\textit{Hint: Treat \(\sin x\) as the variable while factorising, then find every angle in the stated interval.}\par
\(2\sin^2x+\sin x-1=(2\sin x-1)(\sin x+1)\). Hence \(\sin x=\tfrac12\) or \(\sin x=-1\).

In the given interval, \(\sin x=\tfrac12\) gives \(x=\tfrac\pi6\) and \(x=\tfrac{5\pi}6\). The equation \(\sin x=-1\) gives \(x=\tfrac{3\pi}2\).

\textbf{Answer:} \(x=\dfrac\pi6,\ \dfrac{5\pi}6,\ \dfrac{3\pi}2\).

\Needspace{8\baselineskip}\subsection*{2. Parametric differentiation}
\textit{Hint: Use \(\mathrm{d}y/\mathrm{d}x=(\mathrm{d}y/\mathrm{d}t)/(\mathrm{d}x/\mathrm{d}t)\), and find the point as well as the gradient.}\par
At \(t=2\), \(x=5\) and \(y=2\), so the required point is \((5,2)\).

\(\mathrm{d}x/\mathrm{d}t=2t\) and \(\mathrm{d}y/\mathrm{d}t=3t^2-3\). Thus \(\mathrm{d}y/\mathrm{d}x=(3t^2-3)/(2t)\), which is \(9/4\) at \(t=2\).

The tangent through \((5,2)\) with gradient \(9/4\) is \(y-2=\tfrac94(x-5)\).

\textbf{Answer:} \(y-2=\dfrac94(x-5)\).

\Needspace{8\baselineskip}\subsection*{3. Integration by parts}
\textit{Hint: In integration by parts, differentiate \(x\) and integrate \(e^x\).}\par
Take \(u=x\) and \(\mathrm{d}v=e^x\,\mathrm{d}x\), so \(\mathrm{d}u=\mathrm{d}x\) and \(v=e^x\).

\[\begin{aligned}\int_0^1 xe^x\,\mathrm{d}x&=[xe^x]_0^1-\int_0^1e^x\,\mathrm{d}x\\&=e-(e-1)=1.\end{aligned}\]

\textbf{Answer:} \(1\).

\Needspace{9\baselineskip}\subsection*{Find the mistake: correction}
The integral is correctly calculated as zero: the contribution below the axis is \(-1/2\) and the contribution above the axis is \(1/2\).

Total area adds their magnitudes. Equivalently, \(\int_{-1}^1|x|\,\mathrm{d}x=\tfrac12+\tfrac12=1\).

\Needspace{6\baselineskip}\subsection*{What to practise next}\begin{enumerate}
\item Keep the interval and angle units visible when solving trigonometric equations; use a sketch to check for missing roots.
\item Before integrating, decide whether simplification, substitution or integration by parts makes the structure easier.
\item Use graph sketches to check signs, turning points and cancellation before committing to lengthy algebra.
\end{enumerate}\end{document}
